\begin{definition}[$p$-divisible CP map]
    \label{def:p_divisible_channel}
    \lean{MPSTensor.IsPDivisibleChannel}
    \leanok
    A linear endomorphism $\E$ of $\MD$ is \emph{$p$-divisible}
    if there exists a CPTP map $\E'$ of $\MD$ such that
    $\E = (\E')^p$.
\end{definition}

\begin{definition}[$p$-refinable MPS tensor]
    \label{def:p_refinable}
    \lean{MPSTensor.IsPRefinable}
    \leanok
    \uses{def:mps_tensor, def:blocked_tensor}
    Given a tensor $B$, the family $\mc V(B)$ can be \emph{$p$-refined} if
    there exist another tensor $A$ and an isometry
    $W : \C^d \to (\C^d)^{\otimes p}$ such that, for every $N$,
    \begin{align}
        \ket{V^{(pN)}(A)}
         & = W^{\otimes N}\ket{V^{(N)}(B)}.
        \notag
    \end{align}
    Equivalently, after identifying a $p$-letter with an element of
    $\{0,\ldots,d^p-1\}$, the $p$-blocked tensor $A^{[p]}$ satisfies
    \begin{align}
        V^{(N)}(A^{[p]})_\tau
         & =
        \sum_{\sigma \in \{0,\ldots,d-1\}^N}
        \left(\prod_k W_{\tau(k),\sigma(k)}\right)
        V^{(N)}(B)_\sigma,
        \notag
    \end{align}
    for every $N \in \N$ and every
    $\tau \in \{0,\ldots,d^p-1\}^N$.
\end{definition}

\begin{theorem}[The tensor $C$ associated to a refinement witness]
    \label{thm:p_refinement_pullback_irreducible_form}
    \lean{MPSTensor.pRefinementCanonicalization_pullback_of_irreducibleForm}
    \leanok
    \uses{def:irreducible_form, def:p_refinable}
    Let $B$ be an MPS tensor in irreducible form~II. If $\mc V(B)$ can be
    $p$-refined, then there exist a tensor $A$, an isometry
    $W : \C^d \to (\C^d)^{\otimes p}$, and a tensor $C$ defined by
    $C^\alpha := \sum_{i=0}^{d-1} W_{\alpha,i}B^i$ for
    $\alpha \in \{0,\ldots,d^p-1\}$,
    such that $C$ is again in irreducible form~II, $\E_C = \E_B$,
    and $C$ has the same MPV family as $A^{[p]}$.
\end{theorem}

\begin{proof}\leanok
    Choose $(A,W)$ from the refinement data. The definition of $C$ gives
    $\E_C = \E_B$ and equality of MPV families between $C$
    and $A^{[p]}$. To see that $C$ is still in irreducible form~II, transport
    each periodic block $B_j$ of $B$ through the same physical isometry $W$.
    The left-canonical normalization is preserved because
    $W^\dagger W = \Id$,
    while irreducibility and the peripheral spectrum are preserved because the
    transfer map of each transported block agrees with that of $B_j$.
    Reassembling these transported blocks with the same weights gives an
    irreducible-form decomposition of $C$.
\end{proof}

\begin{definition}[$Z$-gauge after applying the blocked equal-case theorem]
    \label{def:peripheral_equalcase_zgauge_sameMPV}
    \lean{MPSTensor.PeripheralEqualCaseZGaugeOfSameMPV}
    \leanok
    \uses{def:irreducible_form, def:blocked_tensor, def:same_mpv, def:z_gauge_equiv}
    Fix $(d,D,p)$. This hypothesis says that whenever a blocked tensor $C$ in
    irreducible form~II has the same MPV family as the $p$-blocked tensor
    $A^{[p]}$,
    there exists an integer $m \ge 1$ and a $\Z_m$-gauge equivalence
    between $C$ and $A^{[p]}$.
\end{definition}

\begin{theorem}[$Z$-gauge equivalence from the periodic theorem]
    \label{thm:zgauge_equiv_periodicFT_irreducible_form}
    \lean{MPSTensor.zGaugeEquiv_of_periodicEqualCaseFT_of_irreducibleForm}
    \leanok
    \uses{thm:periodic_ft, def:irreducible_form, def:blocked_tensor, def:same_mpv, def:z_gauge_equiv}
    Assume the periodic equal-case Fundamental Theorem at the blocked physical
    dimension. Let $C$ and $A^{[p]}$ be MPS tensors of the same bond dimension,
    both in irreducible form~II, and suppose that they have the same MPV family.
    Then there exists an integer $m \ge 1$ and a $\Z_m$-gauge equivalence
    between $C$ and $A^{[p]}$.
    % Scope restriction: this is only the equal-case application of the periodic
    % Fundamental Theorem; the irreducible-form hypothesis on $A^{[p]}$ is an
    % explicit assumption, not a consequence of equality of MPV families.
\end{theorem}

\begin{definition}[Construction of $\widetilde A$ from the $Z$-gauge]
    \label{def:peripheral_equalcase_root_from_zgauge}
    \lean{MPSTensor.PeripheralEqualCaseRootFromZGauge}
    \leanok
    \uses{def:irreducible_form, def:blocked_tensor, def:z_gauge_equiv}
    Fix $(d,D,p)$. This hypothesis says that if $m \ge 1$ and a blocked tensor
    $C$ in irreducible form~II is $\Z_m$-gauge equivalent to $A^{[p]}$,
    then one can construct a tensor $\widetilde A$ with
    $\sum_i (\widetilde A^i)^\dagger \widetilde A^i = \Id$ and
    $\E_C = \E_{\widetilde A^{[p]}}$.
    This channel-level hypothesis is the transfer-map reading of the
    construction in
    \cite[lines~765--810]{DeLasCuevas2017Irreducible}: for each periodic block
    $A_j$, the scalars $c_{j,\alpha}$ define matrices $A'_j$, and then
    $\widetilde A = U^\dagger A'U$.
    This source description is the intended irreducible-form subcase of the
    hypothesis; the definition above keeps the construction as a conditional
    input and does not assert these block scalars $c_{j,\alpha}$ and matrices
    $A'_j$ for arbitrary $A$.
    The cited source moreover obtains the blocked-tensor identity
    $\widetilde A^{[p]} = C$, whereas the formal hypothesis above records only
    its channel-level consequence
    $\E_C = \E_{\widetilde A^{[p]}}$
    (equality of $p$-blocked transfer maps). This conditional input is therefore
    the transfer-map weakening of the source construction, not the full blocked
    tensor equality.
\end{definition}

\begin{definition}[Construction of $\widetilde A$ after the equal-case theorem]
    \label{def:peripheral_equalcase_periodicFT_sameMPV}
    \lean{MPSTensor.PeripheralEqualCasePeriodicFTOfSameMPV}
    \leanok
    \uses{def:irreducible_form, def:blocked_tensor, def:same_mpv}
    Fix $(d,D,p)$. This states the combined construction: whenever a
    blocked tensor $C$ in irreducible form~II has the same MPV family as a
    $p$-blocked tensor $A^{[p]}$, one can replace $A$ by a left-canonical
    tensor $\widetilde A$ with
    $\E_C = \E_{\widetilde A^{[p]}}$.
\end{definition}

\begin{theorem}[Construction of $\widetilde A$ from $Z$-gauge extraction]
    \label{thm:peripheral_equalcase_sameMPV_zgauge_reduction}
    \lean{MPSTensor.peripheralEqualCase_periodicFT_of_sameMPV}
    \leanok
    \uses{def:peripheral_equalcase_zgauge_sameMPV, def:peripheral_equalcase_root_from_zgauge, def:peripheral_equalcase_periodicFT_sameMPV}
    If the blocked equal-case $Z$-gauge extraction hypothesis and the
    construction of $\widetilde A$ from that $Z$-gauge both hold, then the
    combined equal-case construction holds.
\end{theorem}

\begin{proof}\leanok
    Apply the $Z$-gauge extraction hypothesis to obtain a blocked
    $\Z_m$-gauge witness between $C$ and $A^{[p]}$, and then apply the
    construction of $\widetilde A$ to that witness.
\end{proof}

\begin{theorem}[Reduction to the blocked equal-case construction]
    \label{thm:p_refinement_forward_equalcase_reduction}
    \lean{MPSTensor.pRefinementCanonicalization_of_peripheralEqualCase_periodicFT_of_sameMPV}
    \leanok
    \uses{def:peripheral_equalcase_periodicFT_sameMPV, def:irreducible_form, def:p_refinable, thm:p_refinement_pullback_irreducible_form}
    Assume the blocked equal-case construction. Then every
    $p$-refinable tensor $B$ in irreducible form~II admits a left-canonical
    tensor $\widetilde A$ with
    $\E_B = \E_{\widetilde A^{[p]}}$.
\end{theorem}

\begin{proof}\leanok
    Apply Theorem~\ref{thm:p_refinement_pullback_irreducible_form} to the
    refinement witness of $B$. This produces a blocked tensor $C$ in
    irreducible form~II with $\E_C = \E_B$ and the same MPV
    family as $A^{[p]}$. The blocked equal-case construction then
    gives a left-canonical tensor $\widetilde A$ with
    $\E_C = \E_{\widetilde A^{[p]}}$. Combining the two
    transfer-map identities yields
    $\E_B = \E_{\widetilde A^{[p]}}$.
\end{proof}

\begin{theorem}[Forward direction from the blocked equal-case construction]
    \label{thm:thm_4_1_p_refinement_forward_blocked_equalcase}
    \lean{MPSTensor.thm_4_1_p_refinement_forward_of_peripheralEqualCase_periodicFT_of_sameMPV}
    \leanok
    \uses{def:peripheral_equalcase_periodicFT_sameMPV, def:irreducible_form, def:p_refinable, def:p_divisible_channel, thm:p_refinement_forward_equalcase_reduction}
    Assume the blocked equal-case construction. Then for every tensor
    $B$ in irreducible form~II, $p$-refinability of $B$ implies
    $p$-divisibility of $\E_B$.
\end{theorem}

\begin{proof}\leanok
    Combine Theorem~\ref{thm:p_refinement_forward_equalcase_reduction} with the
    witness-level forward implication: once $\E_B$ is identified with
    $\E_{\widetilde A^{[p]}}$ for a left-canonical $\widetilde A$, the
    blocking identity gives $\E_B = (\E_{\widetilde A})^p$.
\end{proof}

The following conditional statements separate the refinability--divisibility
equivalence of \cite[Theorem~4.1]{DeLasCuevas2017Irreducible} from the
additional constructions used below.

% Source comparison: arXiv:1708.00029, Theorem 4.1, lines 728--731 asserts
% the same refinability--divisibility equivalence without additional
% hypotheses. The forward construction of \widetilde A corresponds to
% lines 735--810, while the reverse transfer-map construction corresponds
% to lines 812--818. The Lean theorem below combines the two proved one-way
% implications and is therefore conditional on those two hypotheses.
\begin{theorem}[$p$-refinability iff transfer-map $p$-divisibility, conditional]%
    \label{thm:thm_4_1_p_refinement}
    \lean{MPSTensor.thm_4_1_p_refinement}
    \leanok
    \uses{def:irreducible_form, def:p_refinable, def:p_divisible_channel, thm:periodic_ft, thm:thm_4_1_p_refinement_forward, thm:thm_4_1_p_refinement_reverse}
    Let $B$ be an MPS tensor in irreducible form~II and let $p \ge 1$.
    Assume the forward construction of $\widetilde A$ for
    refinability $\Rightarrow$ divisibility and the reverse transfer-map
    construction for divisibility $\Rightarrow$ refinability. Then $B$ is
    $p$-refinable iff its transfer map $\E_B$ is $p$-divisible.
\end{theorem}

\begin{proof}\leanok
    \uses{thm:thm_4_1_p_refinement_forward, thm:thm_4_1_p_refinement_reverse}
    Combine the forward implication
    (Theorem~\ref{thm:thm_4_1_p_refinement_forward}) with the reverse
    implication (Theorem~\ref{thm:thm_4_1_p_refinement_reverse}) into a
    single biconditional. The forward direction assumes the construction of
    $\widetilde A$; the reverse direction assumes the transfer-map construction.
    Together they give the full equivalence under those hypotheses.
\end{proof}

\begin{theorem}[$p$-refinability implies transfer-map $p$-divisibility]%
    \label{thm:thm_4_1_p_refinement_forward}
    \lean{MPSTensor.thm_4_1_p_refinement_forward}
    \leanok
    \uses{def:irreducible_form, def:p_refinable, def:p_divisible_channel}
    Let $B$ be an MPS tensor in irreducible form~II. Assume that every
    refinement witness gives a tensor $\widetilde A$ with
    $\sum_i (\widetilde A^i)^\dagger \widetilde A^i = \Id$ and
    $\E_B = \E_{\widetilde A^{[p]}}$. Then
    $p$-refinability of $B$ implies $p$-divisibility of $\E_B$.
\end{theorem}
% Formalization status: the assumed construction is refuted, not merely
% unproved (Remark~\ref{rem:thm_4_1_forward_trace_preservation}), so the
% Lean theorem carries an Unfaithful marker and this node stays conditional;
% the elimination plan is in
% docs/paper-gaps/dccsp17_thm41_forward_trace_preservation.tex.

\begin{proof}\leanok
    The construction supplies $\widetilde A$ with
    $\sum_i (\widetilde A^i)^\dagger \widetilde A^i = \Id$ and
    $\E_B = \E_{\widetilde A^{[p]}}$. The left-canonical
    property makes $\E_{\widetilde A}$ a CPTP channel, and the
    blocking identity
    $\E_{\widetilde A^{[p]}} = (\E_{\widetilde A})^p$
    then identifies $\E_B$ with the $p$-th power of the channel
    $\E_{\widetilde A}$.
\end{proof}

\begin{remark}[Trace preservation in the forward direction]
    \label{rem:thm_4_1_forward_trace_preservation}
    \uses{def:irreducible_form, def:p_refinable, def:p_divisible_channel, def:peripheral_equalcase_root_from_zgauge, def:peripheral_equalcase_periodicFT_sameMPV, def:peripheral_equalcase_root_channel_from_zgauge}
    The construction assumed in
    Theorem~\ref{thm:thm_4_1_p_refinement_forward} does not follow from
    irreducible form~II alone, and the printed forward implication of
    \cite[Theorem~4.1]{DeLasCuevas2017Irreducible} omits the trace
    preservation of $\E_B$ that its notion of $p$-divisibility presupposes.
    Let $p \ge 1$, let $B_0$ be $p$-refinable with witness $(A_0, W)$ and in
    irreducible form~II with every multiplicity equal to one, and let
    $\lambda > 1$. Set
    \begin{align}
        B
         & = \lambda B_0,
        \qquad
        A = \lambda^{1/p} A_0.
        \notag
    \end{align}
    Then $B$ is in irreducible form~II with the blocks of $B_0$ and the
    multiplicities $\lambda$, and it is $p$-refinable with witness $(A, W)$:
    \begin{align}
        \ket{V^{(pN)}(A)}
         & = \lambda^{N}\ket{V^{(pN)}(A_0)}
        = \lambda^{N} W^{\otimes N}\ket{V^{(N)}(B_0)}
        = W^{\otimes N}\ket{V^{(N)}(B)}.
        \notag
    \end{align}
    Its transfer map satisfies
    \begin{align}
        \E_B
         & = \lambda^{2}\,\E_{B_0},
        \qquad
        \tr \E_B(\Id) = \lambda^{2}\tr \Id \neq \tr \Id,
        \notag
    \end{align}
    so $\E_B$ is not trace preserving and is not the $p$-th power of a CPTP
    map. At physical and bond dimension one, with $B_0 = (1)$, the same data
    give $C = A^{[p]} = (\lambda)$ related by the trivial $\Z_1$-gauge, and
    no left-canonical $\widetilde A$ or CPTP root has $p$-th power
    $x \mapsto \lambda^2 x$; so the constructions assumed in
    Definitions~\ref{def:peripheral_equalcase_root_from_zgauge},
    \ref{def:peripheral_equalcase_periodicFT_sameMPV},
    and~\ref{def:peripheral_equalcase_root_channel_from_zgauge} fail as well.
    The corrected statement takes $B$ literally in the block form of
    irreducible form~II from \cite{DeLasCuevas2017Irreducible}, as the source
    does, and adds trace preservation of $\E_B$; for such a block form this is
    equivalent to every multiplicity being one, since
    $\sum_i (B^i)^\dagger B^i = \bigoplus_j \mu_j^2 \Id$. Under this
    hypothesis the printed argument closes
    \cite{gap:dccsp17_thm41_forward_trace_preservation}. The equivalence does
    not extend to the generated-family reading of
    Definition~\ref{def:irreducible_form}: a non-unitary similarity of a
    left-canonical normal tensor generates the same family with unit weights,
    while its transfer map is not trace preserving.
\end{remark}

\begin{definition}[Blocked-to-root channel-root hypothesis from $Z$-gauge]
    \label{def:peripheral_equalcase_root_channel_from_zgauge}
    \lean{MPSTensor.PeripheralEqualCaseRootChannelOfZGauge}
    \leanok
    \uses{def:peripheral_equalcase_root_from_zgauge, def:z_gauge_equiv, def:p_divisible_channel}
    Fix $(d,D,p)$. This sharpened reconstruction hypothesis asks for a CPTP map
    $\E'$ such that $\E_C = (\E')^p$ and the Choi/Kraus rank of
    $\E'$ is at most $d$, whenever $C$ is related to $A^{[p]}$ by a
    blocked $\Z_m$-gauge witness.
\end{definition}

\begin{theorem}[Channel-root criterion for blocked-to-root reconstruction]
    \label{thm:peripheral_equalcase_root_from_root_channel}
    \lean{MPSTensor.peripheralEqualCaseRootFromZGauge_of_rootChannel}
    \leanok
    \uses{def:peripheral_equalcase_root_channel_from_zgauge, def:peripheral_equalcase_root_from_zgauge}
    If the blocked-to-root channel-root hypothesis holds, then the tensor-level
    blocked-to-root reconstruction hypothesis holds as well.
\end{theorem}

\begin{proof}\leanok
    Choose a tensor $A'$ whose matrix family realizes the root channel with at
    most $d$ nonzero terms. Since this channel is trace-preserving, the matrices
    of $A'$ satisfy
    $\sum_i (A'^i)^\dagger A'^i = \Id$. The identity
    $\E_C = (\E_{A'})^p$ is then equivalent to
    $\E_C = \E_{A'^{[p]}}$ by the blocking formula.
\end{proof}

\begin{theorem}[Transfer-map $p$-divisibility implies $p$-refinability]%
    \label{thm:thm_4_1_p_refinement_reverse}
    \lean{MPSTensor.thm_4_1_p_refinement_reverse}
    \leanok
    \uses{def:irreducible_form, def:p_refinable, def:p_divisible_channel}
    Let $B$ be an MPS tensor in irreducible form~II and let $p \ge 1$.
    Assume the reverse transfer-map root hypothesis: there is a tensor $A$ with
    $\E_B = \E_{A^{[p]}}$ whenever $\E_B$ is
    $p$-divisible. Then $p$-divisibility of $\E_B$ implies
    $p$-refinability of $B$.
\end{theorem}

\begin{proof}\leanok
    From the reverse transfer-map root hypothesis, extract $A$ with
    $\E_B = \E_{A^{[p]}}$. This matches two finite Kraus
    families of the same completely positive map: the blocked family $A^{[p]}$
    with $d^p$ operators and the original family $B$ with $d$ operators.
    The cardinality comparison $d \le d^p$ holds whenever $p \ge 1$, so
    \cite[Theorem~2.1(4)]{Wolf2012Quantum} supplies an isometry
    $V \in M_{d^p \times d}(\C)$ with $V^\dagger V = \Id$ and
    $(A^{[p]})^\alpha = \sum_j V_{\alpha,j}B^j$. Expanding the
    coefficient of $A^{[p]}$ at the word~$\tau$ by linearity of matrix
    multiplication and of the trace, the contributions reorganise into the
    $V$-weighted sum over length-$N$ words of $B$ required by
    Definition~\ref{def:p_refinable}, exhibiting $V$ as the isometry $W$ in the
    statement of $p$-refinability.
\end{proof}

\begin{remark}[Continuum-limit application]\notready
    \label{rem:thm_4_1_continuum_limit}
    The original motivation for Theorem~\ref{thm:thm_4_1_p_refinement}
    is a continuum-limit construction in the spirit of
    \cite{DeLasCuevas2017Irreducible}: refinability along arbitrarily
    large $p$ is equivalent to infinite divisibility of the transfer
    map, which selects out continuum limits of translationally invariant
    MPS. The further consequences are explored elsewhere.
\end{remark}
